\chapter{Joint Stiffness}

\section*{Definition}
Joint stiffness is the subjective sensation of reduced range of motion or increased resistance to movement of a joint. It is classified as morning stiffness when present upon awakening and lasting more than 30 minutes suggesting inflammatory arthritis, or as gelling stiffness that occurs after periods of inactivity.

\section*{Pathophysiology}
Stiffness in osteoarthritis results from mechanical factors such as cartilage loss, osteophyte formation, and capsular fibrosis limiting smooth gliding. In inflammatory arthritis, synovial effusion, edema, and pannus formation cause mechanical obstruction, while cytokines (TNF-$\alpha$, IL-6) increase joint fluid viscosity. Non-articular stiffness may arise from myofascial restriction, tendon shortening, or fibrosis of periarticular structures.

\section*{Causes}
\begin{itemize}
  \item \textbf{Inflammatory arthritis:} Rheumatoid arthritis, psoriatic arthritis, ankylosing spondylitis, systemic lupus erythematosus
  \item \textbf{Degenerative:} Osteoarthritis (typically brief morning stiffness $<$ 30 minutes)
  \item \textbf{Infectious:} Septic arthritis, Lyme disease, viral polyarthritis
  \item \textbf{Crystal arthropathy:} Gout, pseudogout
  \item \textbf{Metabolic:} Hemochromatosis, hypothyroidism, diabetes mellitus (cheiroarthropathy)
  \item \textbf{Medication-related:} Statins, aromatase inhibitors, fluoroquinolones
  \item \textbf{Myofascial:} Fibromyalgia, polymyalgia rheumatica, myofascial pain syndrome
\end{itemize}

\section*{When to See a Doctor}
See your doctor if:
\begin{itemize}
  \item Morning stiffness lasting longer than 30--60 minutes
  \item Joint stiffness with swelling, warmth, or erythema
  \item Stiffness impairing activities of daily living
  \item Associated systemic symptoms: fever, weight loss, fatigue, rash
\end{itemize}

\section*{Related Symptoms}
\begin{itemize}
  \item Joint pain, tenderness, and swelling
  \item Reduced range of motion and functional limitation
  \item Warmth and erythema over the affected joint
  \item Systemic symptoms (fever, malaise, weight loss) in inflammatory disease
  \item Symmetric vs. asymmetric joint involvement pattern
\end{itemize}
\textbf{Important:} Prolonged morning stiffness (> 60 minutes) is a hallmark of inflammatory arthritis and should prompt serologic evaluation for rheumatoid arthritis or spondyloarthropathy.

\section*{Self-Care / Home Management}
\begin{itemize}
  \item Gentle range-of-motion exercises and stretching upon waking
  \item Warm showers or moist heat packs applied to stiff joints for 15--20 minutes
  \item Paraffin wax baths for hands and feet
  \item Over-the-counter NSAIDs or acetaminophen for symptomatic relief
  \item Ergonomic modifications to reduce joint stress during daily tasks
\end{itemize}

\section*{How Your Doctor Will Evaluate You}
\begin{itemize}
  \item Detailed history: duration, pattern, associated symptoms, family history
  \item Physical examination assessing range of motion, swelling, tenderness, and extra-articular signs
  \item Laboratory evaluation: ESR, CRP, RF, anti-CCP, ANA, uric acid, TSH, CPK
  \item Plain radiography of affected joints for erosions, joint space narrowing, osteophytes
  \item Ultrasound or MRI to evaluate synovitis, effusion, and soft tissue involvement
\end{itemize}

\section*{Treatment Options}
\begin{itemize}
  \item NSAIDs or cyclooxygenase-2 (COX-2) inhibitors for symptom relief
  \item Disease-modifying antirheumatic drugs (methotrexate, sulfasalazine, hydroxychloroquine) for inflammatory arthritis
  \item Biologic therapy (TNF inhibitors, IL-6 inhibitors) for refractory inflammatory disease
  \item Corticosteroid injections for monoarticular flares
  \item Physical and occupational therapy for joint preservation and mobility
  \item Weight loss and low-impact aerobic exercise for osteoarthritis
\end{itemize}

\clearpage
